\documentclass[10pt]{article}

\usepackage{amsmath}
\usepackage{amssymb}
\usepackage{amsthm}
\usepackage{ mathrsfs }
\usepackage{ mathtools}
\usepackage{enumerate}
\usepackage{bbm}
\usepackage{lipsum}
\usepackage{fancyhdr}
\usepackage{tikz-cd} 
\usepackage{adjustbox} 
\usetikzlibrary{arrows}
\newcommand{\midarrow}{\tikz \draw[-triangle 90] (0,0) -- +(.1,0);}
\tikzset{commutative diagrams/.cd,
mysymbol/.style={start anchor=center,end anchor=center,draw=none}
}
\newcommand\comsymb[2][\circlearrowleft]{%
  \arrow[mysymbol]{#2}[description]{#1}}

\newtheorem{theorem}{Theorem} [section] 
\newtheorem{proposition}{Proposition}[section] 
\newtheorem{definition}{Definition}[section] 
\newtheorem{lemma}{Lemma}[section] 
\newtheorem{notation}{Notation}[section] 
\newtheorem{remark}{Remark}[section] 
\newtheorem{corollary}{Corollary} [section] 
\newtheorem{terminology}{Terminology}[section] 
\newtheorem{fact}{Fact}[section] 
\newtheorem{example}{Example}[section] 
\newtheorem{claim}{Claim}
\newenvironment{claimproof}[1]{\par\noindent\underline{Proof:}\space#1}{\hfill $\blacksquare$}

\DeclareMathOperator{\tmod}{mod}
\DeclareMathOperator{\triv}{triv}
\DeclareMathOperator{\ind}{ind}
\DeclareMathOperator{\straw}{s.t.}
\DeclareMathOperator{\ev}{ev}
\DeclareMathOperator{\pr}{pr}
\DeclareMathOperator{\Aut}{Aut}
\DeclareMathOperator{\Map}{Map}
\DeclareMathOperator{\Stab}{Stab}
\DeclareMathOperator{\Ind}{Ind}
\DeclareMathOperator{\GL}{GL}
\DeclareMathOperator{\SL}{SL}
\DeclareMathOperator{\SO}{SO}
\DeclareMathOperator{\Sp}{Sp}
\DeclareMathOperator{\esssup}{esssup}
\DeclareMathOperator{\abs}{abs}
\DeclareMathOperator{\rel}{rel}
\DeclareMathOperator{\diam}{diam}
\DeclareMathOperator{\rank}{rank}
\DeclareMathOperator{\Hom}{Hom}
\DeclareMathOperator{\Homk}{Hom_{k-alg}}
\DeclareMathOperator{\Ker}{Ker}
\DeclareMathOperator{\Image}{Im}
\DeclareMathOperator{\Dom}{Dom}
\DeclareMathOperator{\grad}{grad}
\DeclareMathOperator{\divergence}{div}
\DeclareMathOperator{\rk}{rk}
\DeclareMathOperator{\Span}{Span}
\DeclareMathOperator{\mSpec}{mSpec}
\DeclareMathOperator{\MaxSpec}{MaxSpec}
\DeclareMathOperator{\interior}{int}
\DeclareMathOperator{\supp}{supp}
\DeclareMathOperator{\id}{id}
\DeclareMathOperator{\sgn}{sgn}
\DeclareMathOperator{\Mat}{Mat}
\DeclareMathOperator{\PD}{PD}
\DeclareMathOperator{\ext}{ext}
\DeclareMathOperator{\vol}{vol}
\DeclareMathOperator{\inte}{int}
\DeclareMathOperator{\diverg}{div}
\DeclareMathOperator{\curl}{curl}
\DeclareMathOperator{\image}{im}
\DeclareMathOperator{\dR}{dR}
\DeclareMathOperator{\Ob}{Ob}
\DeclareMathOperator{\Mnfd}{Mnfd}
\DeclareMathOperator{\OpEmb}{OpEmb}
\DeclareMathOperator{\Spec}{Spec}
\DeclareMathOperator{\Spa}{Spa}
%\newcommand*{\name}[\num_arguments][default values]{{\color{#1}\Large #2}}
\newcommand{\defeq}{\vcentcolon=}
\newcommand{\Ouv}{\mathcal{O}}
\newcommand{\norm}[1]{\Vert #1 \Vert}
\newcommand{\opNorm}[2]{\norm{#1}_{#2\to#2}}
\newcommand{\normL}[3]{\norm{#1}_{L^{#2}(#3)}}
\newcommand{\N}[0]{\mathbb{N}}
\newcommand{\m}[0]{\mathfrak{m}}
\newcommand{\R}[0]{\mathbb{R}}
\newcommand{\Z}[0]{\mathbb{Z}}
\newcommand{\C}[0]{\mathbb{C}}
\newcommand{\Q}[0]{\mathbb{Q}}
\newcommand{\F}[0]{\mathbb{F}}
\newcommand{\G}[0]{\mathbb{G}}
\newcommand{\A}[0]{\mathbb{A}}
\newcommand{\T}[0]{\mathbb{T}}
\newcommand{\Para}[0]{\mathbb{P}}
\newcommand{\M}[0]{\mathcal{M}}
\newcommand{\maniN}[0]{\mathcal{N}}
\newcommand{\cinf}[0]{\mathcal{C}^\infty}
\newcommand{\torus}[0]{\mathbb{T}}
\newcommand{\sep}[0]{\:|\:}
\newcommand{\pd}[2]{{\frac {\partial #1} {\partial #2}}}
\newcommand{\compinc}[0]{\subset \subset}
\newcommand{\sH}[0]{\mathscr{H}}
\newcommand{\ab}[1]{\vert #1 \vert}
\newcommand{\st}[0]{\:\straw\:}
\newcommand{\fib}[1]{%
  \mathbin{\mathop{\times}\limits_{#1}}%
}
\newcommand{\tens}[1]{%
  \mathbin{\mathop{\otimes}\displaylimits_{#1}}%
}

\title{Rigid Analytic Geometry Week 1 Solution}
\author{So Murata}
\date{SoSe 26, University of Bonn}


\begin{document}
\maketitle

\par\textbf{Problem 1}\\
By multiplicative property of the absolute value, we have $K^\circ$ is a valuation ring of $K$. Thus $\Spec K^\circ = \{K^{\circ\circ},\{0\}\}$.
\par\textbf{Problem 2}\\
The second assertion is clear as for any $\alpha\in\Z_{\geq0}^m$, $1_\alpha=0$ unless $\alpha = 0$. For the first assertion, from the definition of $\norm{\cdot|\T_m}$, for each $f,g\in\T_m$ there are $\alpha,\beta\in\Z_{\geq0}^m$ such that 
\begin{equation}
    \vert f_\alpha\vert=\norm{f|\T_m},\vert g_\beta\vert = \norm{g|\T_m}.
\end{equation}
Also,
\begin{align*}
\norm{fg|\T_m} &= \max_\gamma\left\{\left\vert\sum_{\alpha+\beta = \gamma}f_\alpha g_\beta\right\vert\sep \gamma\in\Z_{\geq0}\right\},\\
&\leq\vert f_\alpha\vert\vert g_\beta\vert,\\
& = \norm{f|\T_m}\norm{g|\T_m}.
\end{align*}
To show multiplicativeness, consider,
\begin{equation*}
  \alpha_0 = \min\{\alpha\sep \vert f_\alpha\vert = \norm{f|\T_m}\},  \beta_0 = \min\{\beta\sep \vert g_\beta\vert = \norm{g|\T_m}\}.
\end{equation*}

Then,
\begin{equation*}
  \vert (fg)_{\alpha_0+\beta_0}\vert \leq \norm{fg|\T_m}.
\end{equation*}

For any $(\alpha,\beta)$ with $\alpha+\beta = \alpha_0+\beta_0$ we have,
\begin{equation*}
  \vert f_\alpha g_\beta\vert <\vert f_{\alpha_0}g_{\beta_0}\vert.
\end{equation*}

As $\vert\cdot\vert$ is non-archimedian, we have,
\begin{equation*}
  \left\vert\sum_{\substack{\alpha+\beta=\alpha_0+\beta_0\\(\alpha,\beta)\not=(\alpha_0,\beta_0)}}f_\alpha g_\beta\right\vert <\vert f_{\alpha_0}g_{\beta_0}\vert.
\end{equation*}

This can be seen by induction. Using it is non-archimedian, we have,
\begin{align*}
  \vert f_{\alpha_0}g_{\beta_0}\vert&=\left\vert (fg)_{\alpha_0+\beta_0}-\sum_{\substack{\alpha+\beta=\alpha_0+\beta_0\\(\alpha,\beta)\not=(\alpha_0,\beta_0)}}f_\alpha g_\beta\right\vert,\\
  &\leq \max\left\{\vert (fg)_{\alpha_0+\beta_0}\vert,\left\vert\sum_{\substack{\alpha+\beta=\alpha_0+\beta_0\\(\alpha,\beta)\not=(\alpha_0,\beta_0)}}f_\alpha g_\beta\right\vert \right\}
\end{align*}

But 
\begin{equation*}
    \left\vert\sum_{\substack{\alpha+\beta=\alpha_0+\beta_0\\(\alpha,\beta)\not=(\alpha_0,\beta_0)}}f_\alpha g_\beta\right\vert <\vert f_{\alpha_0}g_{\beta_0}\vert.
\end{equation*}

Thus we conclude
\begin{equation*}
  \vert (fg)_{\alpha_0+\beta_0}\vert = \vert f_{\alpha_0}g_{\beta_0}\vert,
\end{equation*}
showing this is multiplicative.

\par\textbf{Probmel 3}\\
Clearly $0,1$ are idempotents thus $0,1\in A^\circ$. From the first property of the norm, 
\begin{equation*}
\norm{ab|A}\leq\norm{a|A}\norm{b|A},
\end{equation*}
$\{(ab)^n\sep n\in\N\}$ is bounded by the products of upperbounds of $\{a^n\sep n\in\N\},\{b^n\sep n\in\N\}$, thus $ab\in A^\circ$. Similarly for $a+b$, as
\begin{equation*}
    \norm{(a+b)^n|A}\leq \norm{a+b|A}^n\leq \max\{\norm{a|A},\norm{b|A}\}^n.
\end{equation*}
Thus we can take the maximum of the upperbounds of $\{a^n\sep n\in\N\},\{b^n\sep n\in\N\}$.
\par If $a\in A^\circ$ and $x\in A^{\circ\circ}$ then as $x^n\stackrel{n\to\infty}{\to}0$ and there is $M>0$ such that $\norm{a|A}^n\leq M$, we have,
\begin{equation*}
    \norm{(ax)^n|A}\stackrel{n\to\infty}{\to}0.
\end{equation*}
Since $A$ is a complete metric space with respect to the induced metric from $\norm{\cdot|A}$, $0$ is the unique element of norm $0$, thus $ax\in A^{\circ\circ}$. If $a^k\in A^{\circ\circ}$, then again by that $A$ is a complete metric space we have $a\in A^{\circ\circ}$ since $\norm{a^{kn}|A}\stackrel{n\to\infty}{\to}0$.\\
\par\textbf{Problem 4}\\
If $\norm{\cdot|A}$ is multiplicative, the boundedness of the set $\{a^n\sep n\in\N\}$ is equivalent to the boundedness of $\{\norm{a|A}^n\sep n\in\N\}$. And the latter is equivalent to $\norm{a|A}\leq 1$. Similarly $\lim_{n\to\infty}a^n=0$ is equivalent to $\norm{a^n|A}\to 0$ and by the multiplicativeness this is equivalent to $\lim_{n\to\infty}\norm{a|A}^n=0$. This is equivalent to $\norm{a|A}<1$.
\\\par\textbf{Problem 5}\\
Take $a\in K$ such that $\vert a \vert\in(0,1)$. Then for a $K$-Banach algebra $A$, we have $a\cdot 1_A$ is a topologically nilpotent element.
As any ideal $I$ contains $0$, thus if $I$ is open, it cotains an open neighbourhood of $0$. This means that for sufficiently large $n$ we have $(a\cdot 1_A)^n\in I$. Thus $1_A\in I$ meaning $I$ is not proper.
\\\par\textbf{Problem 6}
Let $A$ be a $K$-Banach algebra. Clearly $1_A\in A^\times$. If $A^\times$ contains an open neighbourhood $U$ of $1$, then for any $a\in A^\times, aU$ is an open neighbourhood of $a$ contained in $A^\times$. We set,
\begin{equation*}
  U = \{x\in A\sep \norm{x|A}<1\}.
\end{equation*}
Clearly $U$ is open. Note that for any $a\in A$, $x\mapsto a+x$ is a continuous function with inverse $x\mapsto -a+x$. And as $A$ is a $K$-Banach algebra, this is a bounded linear operator. Thus $1+U$ is homeomorphic to $U$.
Consider,
\begin{equation*}
  (1-x)(1+x+x^2+\cdots+x^{n-1}) = 1-x^{n-1}.
\end{equation*}
As $\norm{x|A}<1$, thus the series, 
\begin{equation*}
  1+x+x^2+\cdots
\end{equation*}
is an element of $A$, we conclude $1-x\in A^\times$. In other words, $1+U$ is an open set contained in $A^\times$. Thus $A^\times$ is open.
\\\par\textbf{Problem 7}
By Problem 6, we know that in a $K$-Banach algebra $A$, $A^\times$ forms an open neighbourhood of $1$. If an ideal $I$ in $A$ is dense, $I$ intersects with $A^\times$. This means $I$ contains a unit thus $I=A$.
\end{document}